\originalsection{18.112第二次历史考试（2006）}
\sourceinfo{题面mid2prob.pdf, PDF p.2实印“Nov. 22, 2006”, 开卷; 官方解答mid2.pdf, pp.2--4. MIT OCW 18.112 Fall 2008归档, CC BY-NC-SA 4.0. 4题, 各分值保留.}
\begin{exercise}\label{ass:m2-1}
\textbf{原题 M2.1（15分）。} $\gamma$ 为正向圆 $|z|=9$, 求 $\int_\gamma(\ee^z-1)^{-1}\dd z$.
\end{exercise}
\begin{solution}据官方解答翻译. 全部极点 $2k\pi\ii$ 均一阶, 留数1. 圆内恰有 $k=0,1,-1$ 三点, 圆上无极点. 留数定理给 $6\pi\ii$.\end{solution}
\begin{exercise}\label{ass:m2-2}
\textbf{原题 M2.2（30分）。} 整函数 $f$ 满足 $(\Re f(z))/z\to0$ 于 $z\to\infty$, 证明 $f$ 常数. 提示使用 $|z|<R$ 时的 Schwarz 公式
\[
f(z)=\frac1{2\pi\ii}\int_{|\zeta|=R}u(\zeta)\frac{\zeta+z}{\zeta-z}\frac{\dd\zeta}{\zeta}+\ii C,\qquad u=\Re f,
\]
并仔细估计 $|z|<R/2$ 时 $f'(z)$. 原卷把 $u=\Re f$ 一处误印为 $u=\Re(x)$, 本式已注出应读的意义. 条件的分子是实部, 不默改为 $\Re(f(z)/z)$.
\end{exercise}
\begin{solution}
据官方解答翻译并补足极限量词; 直接令估计趋0的收尾为编者补充（AI）. 条件等价于 $|u(z)|/|z|\to0$. 固定任意 $z$, 对每个 $\varepsilon>0$ 取充分大的 $R>2|z|$, 使圆周 $|u(\zeta)|\le\varepsilon R$. Schwarz 公式在紧子圆盘可对 $z$ 求导, 因为分母与积分路径分离, 得 $f'(z)=(\pi\ii)^{-1}\int_{|\zeta|=R}u(\zeta)/(\zeta-z)^2\dd\zeta$. 从而
\[
|f'(z)|\le\frac1\pi\frac{2\pi R\,\varepsilon R}{(R-|z|)^2}\le8\varepsilon.
\]
任意小 $\varepsilon$ 均成立, 故 $f'(z)=0$. 任意 $z$ 都可如此处理, $f$ 常数. 原解先取上界8, 对整函数 $f'$ 用 Liouville 得线性函数, 再排除斜率; 此处的 $\varepsilon$ 估计直接完成同一结论.
\end{solution}
\begin{exercise}\label{ass:m2-3}
\textbf{原题 M2.3（25分）。} $f$ 在 $|z|<1$ 解析且 $|f(z)|\le(1-|z|)^{-1}$, 求 Cauchy 公式能够给出的 $f^{(n)}(0)$ 最佳估计. 提示在每个 $|z|\le r<1$ 使用公式.
\end{exercise}
\begin{solution}
据官方解答翻译并补出 $n=0$ 的情形. 任意 $0<r<1$ 给 $|f^{(n)}(0)|\le n!/((1-r)r^n)$. 对整数 $n\ge1$, 最大化 $(1-r)r^n$: 其对数导数为 $n/r-1/(1-r)$, 唯一最大点 $r=n/(n+1)$; 两端乘积趋0. 最佳圆周上界为 $n!(n+1)^{n+1}/n^n=(n+1)!(1+1/n)^n$. $n=0$ 时取 $r\downarrow0$, 给 $|f(0)|\le1$, 也可直接代入条件. “Cauchy 公式给出的最佳”指在此给定模界下优化半径, 不在本题额外声称对所有函数的最优极值都已实现.
\end{solution}
\begin{exercise}\label{ass:m2-4}
\textbf{原题 M2.4（30分）。} $P(z)=z^7-2z^5+6z^3-z+1$ 在 $|z|<1$ 中有多少根? 在 $|z|=2$ 所围内部有多少根? 提示在单位圆找最大项并用 Rouché 定理.
\end{exercise}
\begin{solution}据官方解答翻译, 根均计重数. 单位圆上 $|P-6z^3|\le1+2+1+1=5<6$, 故内部3根. 半径2圆上 $|P-z^7|\le64+48+2+1=115<128=|z^7|$, 故内部7根. 两圆边界均无根, 所以未遗漏边界情形.\end{solution}
